\section{Introduction}
\label{sec:introduction}
Understanding the causal effect of a continuous variable (termed ``treatment'') on individual units and subgroups is crucial across many fields. 
In economics, we might like to know the effect of price on demand from different customer demographics. 
In healthcare, we might like to know the effect of medication dosage on health outcomes for patients of various ages and comorbidities. 
And in climate science, we might like to know the effects of anthropogenic emissions on cloud formation and lifetimes under variable atmospheric conditions. 
In many cases, these effects must be estimated from observational data as experiments are often costly, unethical, or otherwise impossible to conduct. 

Estimating causal effects from observational data can only be done under certain conditions, some of which are not testable from data. 
The most prominent are the common assumptions that all confounders between treatment and outcome are measured (``no hidden confounders''), and any level of treatment could occur for any observable covariate vector (``positivity''). 
These assumptions and their possible violations introduce uncertainty when estimating treatment effects. 
Estimating this uncertainty is crucial for decision-making and scientific understanding.
For example, understanding how unmeasured confounding can change estimates about the impact of emissions on cloud properties can help to modify global warming projection models to account for the uncertainty it induces.  

We present a novel marginal sensitivity model for continuous treatment effects. 
This model is used to develop a method that gives the user a corresponding interval representing the ``ignorance region'' of the possible treatment outcomes per covariate and treatment level \citep{d2019multi} for a specified level of violation of the no-hidden confounding assumption. 
We adapt prior work \citep{tan2006msm,kallus2019interval,jesson2021quantifying} to the technical challenge presented by continuous treatments. 
Specifically, we modify the existing model to work with propensity score densities instead of propensity score probabilities (see \Cref{sec:cmsm} below) and propose a method to relate ignorability violations to the unexplained range of outcomes.
% \uri{should add that we add a method for ground the MSM violation in terms of outcomes (the $\rho$ thing)}
Further, we derive bootstrapped uncertainty intervals for the estimated ignorance regions and show how to efficiently compute the intervals, thus providing a method for quantifying the uncertainty presented by finite data and possible violations of the positivity assumption. 
We validate our methods on synthetic data and provide an application on real-world satellite observations of the effects of anthropogenic emissions on cloud properties. 
For this application, we develop a new neural network architecture for estimating continuous treatment effects that can take into account spatiotemporal covariates. 
We find that the model accurately captures known patterns of cloud deepening in response to anthropogenic emission loading with realistic intervals of uncertainty due to unmodeled confounders in the satellite data.
% \uri{we might want to mention here the ``grounding'' idea we discussed on Monday}.